Bei konstanter Sinkgeschwindigkeit ist die resultierende Kraft null: Die Gewichtskraft ist gleich gross wie die Summe aus Auftriebskraft und Strömungswiderstand. Für eine Kugel mit Radius $r$ (Volumen $\frac{4}{3}\pi r^3$, Stirnfläche $\pi r^2$) gilt mit den Dichten $\rho_\text{S}$ von Stahl und $\rho_\text{W}$ von Wasser
\begin{align}
 \rho_\text{S}\cdot\frac{4}{3}\pi r^3\,g &= \rho_\text{W}\cdot\frac{4}{3}\pi r^3\,g + \frac{1}{2}\,c_w\,\rho_\text{W}\,\pi r^2\,v^2
\end{align}
Aufgelöst nach $v$:
\begin{align}
 v &= \sqrt{\frac{8\left(\rho_\text{S}-\rho_\text{W}\right)g}{3\,c_w\,\rho_\text{W}}\cdot r} \quad\Rightarrow\quad v \propto \sqrt{r}
\end{align}
Da $c_w$ und die Dichten für beide Kugeln gleich sind, ist
\begin{align}
 \frac{v_1}{v_2} &= \kVeF = \sqrt{\kRaO} = \kVe \approx \result{\kVeP{0}{2}}
\end{align}
Die grössere Kugel sinkt also um den Faktor $\sqrt{2}$ schneller.
